\title{DeepSeekMath: Pushing the Limits of Mathematical Reasoning in Open Language Models}

\begin{abstract}
Mathematical reasoning remains a particularly demanding task for language models, stemming from its inherent structure and complexity. This work presents DeepSeekMath~7B, developed by further pre-training DeepSeek-Coder-Base-v1.5 7B on a corpus containing 120B math-specific tokens collected from Common Crawl, integrated with both natural language and programming data. DeepSeekMath~7B attains a remarkable 51.7\% accuracy on the MATH benchmark, a competitive examination-level test, without employing external tools or ensemble voting strategies—achieving performance close to that of Gemini-Ultra and GPT-4. Furthermore, by sampling 64 outputs and applying self-consistency, the score increases to 60.9\% on the MATH dataset. DeepSeekMath~’s enhanced mathematical reasoning abilities arise from two principal elements: Firstly, the model exploits the vast reservoir of freely available online data through a rigorously crafted data selection process. Secondly, it utilizes Group Relative Policy Optimization (GRPO), an adaptation of Proximal Policy Optimization (PPO), to simultaneously improve reasoning proficiency in mathematics and reduce PPO’s memory demands.
\end{abstract}

\section{Introduction}

Large language models (LLMs) have transformed mathematical reasoning within artificial intelligence, leading to remarkable progress on benchmarks designed for both quantitative reasoning \citep{MATH} and geometric reasoning \citep{trinh2024solving}. In addition, LLMs have emerged as valuable tools for supporting humans in resolving intricate mathematical challenges \citep{tao}. Nevertheless, state-of-the-art systems like GPT-4 \citep{gpt4} and Gemini-Ultra \citep{gemini} remain closed-source, while available open-source alternatives continue to lag substantially in their mathematical reasoning capabilities.

In this paper, we present DeepSeekMath, a language model specialized for the mathematics domain that demonstrates notable superiority over other open-source models in mathematical reasoning and nearly reaches GPT-4’s benchmark results in academic settings. This advancement is made possible through the development of the DeepSeekMath Corpus—a comprehensive, large-scale pre-training resource containing 120B math-specific tokens. The dataset was curated from Common Crawl (CC) data by leveraging a classifier based on fastText \citep{joulin2016fasttext}. Initially, this classifier was trained using examples drawn from OpenWebMath \citep{openwebmath} as positives, with a broad array of unrelated web content functioning as negative samples. Using this preliminary model, we extracted further candidate positive samples from the CC corpus; these selections underwent manual review and annotation. Subsequently, the classifier was retrained on this enriched dataset, yielding improved extraction precision. Our empirical assessments reveal that this large-scale resource is of substantial quality: the DeepSeekMath-Base 7B model attains a score of 64.2\% on GSM8K \citep{gsm8k} and 36.2\% on the MATH competition dataset \citep{MATH}, surpassing Minerva 540B \citep{minerva} in performance. Furthermore, owing to the multilingual nature of the DeepSeekMath Corpus, we observe enhanced outcomes on Chinese mathematical evaluation sets \citep{wei2023cmath,agieval}. We suggest that our methodology for curating mathematical data offers a useful foundation for the research community and anticipate considerable opportunities for progress in this field.

DeepSeekMath-Base is constructed by initializing from DeepSeek-Coder-Base-v1.5 7B \citep{deepseek-coder}, as empirical evidence suggests that a model pre-trained on code yields superior outcomes relative to initializing with a general-purpose language model. Notably, mathematical pre-training not only strengthens mathematical proficiency but also enhances the model's general reasoning, as reflected by improvements on MMLU \citep{mmlu} and BBH \citep{bbh} evaluations.

Following this pre-training phase, mathematical instruction tuning is performed on DeepSeekMath-Base, utilizing data sources with chain-of-thought \citep{cot}, program-of-thought \citep{pot,pal}, and tool-augmented reasoning \citep{tora} annotations. The tuned model, DeepSeekMath-Instruct 7B, surpasses all other models of comparable scale (7B), and attains performance comparable to instruction-tuned models with 70B parameters in the open-source community.

In addition, we propose Group Relative Policy Optimization (GRPO), an alternative reinforcement learning approach that adapts Proximal Policy Optimization (PPO) \citep{schulman2017proximal}. GRPO eliminates the need for a critic network and instead derives the baseline through group-level scoring, leading to significant savings in training computation. By employing only a selected portion of English instruction-tuning samples, GRPO yields marked gains beyond the already robust DeepSeekMath-Instruct. These advances manifest across both in-domain benchmarks (GSM8K: 82.9\% $\rightarrow$ 88.2\%, MATH: 46.8\% $\rightarrow$ 51.7\%) and out-of-domain problems (e.g., CMATH: 84.6\% $\rightarrow$ 88.8\%) during the RL stage. We further introduce a unifying framework to interpret several methods, encompassing Rejection Sampling Fine-Tuning (RFT) \citep{yuan2023scaling}, Direct Preference Optimization (DPO) \citep{dpo}, PPO, and GRPO, showing that these can all be viewed as variants of direct or streamlined RL procedures. Extensive analyses are conducted to dissect the critical factors of this framework, considering distinctions such as online versus offline training, supervision on results versus processes, and single-step versus iterative reinforcement learning. Finally, we articulate how our RL approach enhances the abilities of instruction-tuned models, and outline avenues for further refining RL methodologies within the proposed unified framework.

\subsection{Contributions}
We present advances in scalable mathematical pre-training and conduct a systematic study of reinforcement learning methodologies.

\textbf{Math Pre-Training at Scale}

\begin{itemize}[topsep=0pt]
    \item Through our investigation, we demonstrate that the widely available Common Crawl dataset harbors extensive mathematical information. Employing a carefully crafted data filtering pipeline, we create the DeepSeekMath Corpus—a rigorously curated collection comprising 120B tokens drawn from web pages with mathematical content. This dataset is nearly 7 times larger than that used by Minerva \citep{minerva} and about 9 times bigger than the recently published OpenWebMath \citep{openwebmath}.

    \item DeepSeekMath-Base 7B, our base pre-trained model, attains performance on par with Minerva 540B \citep{minerva}. This suggests that parameter count alone does not dictate a model’s proficiency in mathematical reasoning; instead, models of modest size, when trained on high-quality datasets, can demonstrate strong mathematical abilities.

    \item We report insights from our pre-training experiments. Performing code training before mathematics training significantly enhances a model's mathematical problem-solving capacities, both with and without auxiliary tools. This partially addresses a persistent question in the field: \textit{does training on code enhance reasoning skills?} Our results suggest this is indeed the case, particularly for mathematical reasoning.

    \item Although utilizing arXiv as a training source is widespread—especially for math-centric research—we observe no meaningful gains on any of the mathematical benchmarks evaluated in this work.
\end{itemize}

\textbf{Exploration and Analysis of Reinforcement Learning}

\begin{itemize}[topsep=0pt]
    \item This work presents Group Relative Policy Optimization (GRPO), a novel reinforcement learning technique characterized by its efficiency and effectiveness. GRPO eliminates the reliance on a critic model and instead leverages group scores to compute the baseline, thereby markedly reducing the computational overhead during training in comparison to Proximal Policy Optimization (PPO).
    \item Our empirical analyses reveal that GRPO markedly boosts the capabilities of DeepSeekMath-Instruct, our instruction-tuned model, by training exclusively with instruction-tuning datasets. Additionally, we report notable improvements in generalization to out-of-domain tasks observed throughout the reinforcement learning process.
    \item We establish a cohesive conceptual framework that encapsulates a range of approaches, including RFT, DPO, PPO, and GRPO. Through comprehensive experimentation—encompassing online versus offline training, outcome versus process supervision, and single-turn versus iterative reinforcement learning—we thoroughly examine the core components of this unified approach.
    \item Leveraging insights from this unified perspective, we investigate the mechanisms that underlie the success of reinforcement learning, and articulate several promising avenues for enhancing the reinforcement learning of LLMs in future research.
\end{itemize}

\subsection{Summary of Evaluations and Metrics}

\begin{itemize}[topsep=0pt]
    \item \textbf{English and Chinese Mathematical Reasoning}: We systematically evaluate our models against a variety of benchmarks in both English and Chinese, which span mathematical reasoning tasks from elementary school up to university curricula. The English benchmarks considered comprise GSM8K \citep{gsm8k}, MATH \citep{MATH}, SAT \citep{llemma}, OCW Courses \citep{minerva}, and MMLU-STEM \citep{mmlu}. For Chinese, we include MGSM-zh \citep{mgsm}, CMATH \citep{wei2023cmath}, Gaokao-MathCloze \citep{agieval}, and Gaokao-MathQA \citep{agieval}. Our evaluations cover models’ capacities to produce fully self-explanatory text-based solutions (in the absence of auxiliary tools) as well as to solve problems programmatically with Python.

    On the suite of English benchmarks, DeepSeekMath-Base exhibits performance competitive with the closed-source Minerva 540B \citep{minerva}, and outperforms all current open-source base models (including Mistral 7B \citep{mistral} and Llemma-34B \citep{llemma}), regardless of their exposure to mathematical pre-training, frequently by considerable margins. Strikingly, DeepSeekMath-Base also achieves superior results on Chinese tasks, which we attribute to our deviation from prior approaches \citep{minerva,llemma} that restrict pre-training data to English, opting instead to include high-caliber datasets in other languages. Upon further instruction tuning and reinforcement learning, DeepSeekMath-Instruct and DeepSeekMath-RL attain new levels of performance, notably becoming the first open-source models to exceed 50\% accuracy on the competitive MATH benchmark.
\end{itemize}

\item \textbf{Formal Mathematics}: We assess DeepSeekMath-Base on the informal-to-formal theorem proving benchmark introduced in \citep{dsp_proof} using miniF2F \citep{minif2f}, where Isabelle \citep{isabelle} serves as the proof assistant. DeepSeekMath-Base exhibits notable proficiency in few-shot autoformalization for this task.
\item \textbf{Natural Language Understanding, Reasoning, and Code}: For a holistic evaluation of the model’s abilities across general comprehension, reasoning, and coding, we test DeepSeekMath-Base on several benchmarks: the Massive Multitask Language Understanding (MMLU) benchmark \citep{mmlu}, which features 57 multiple-choice tasks spanning a broad array of topics; BIG-Bench Hard (BBH) \citep{bbh}, which poses 23 especially demanding tasks, most requiring complex, multi-step reasoning; as well as HumanEval \citep{codex} and MBPP \citep{mbpp}, both of which are established measures for code model performance. Results indicate that mathematical pre-training confers gains not only for mathematical reasoning but also for language comprehension and general reasoning tasks.
\end{itemize}

\section{Math Pre-Training}

\subsection{Data Collection and Decontamination} 
Here, we describe how the DeepSeekMath Corpus is constructed from the Common Crawl dataset. Illustrated in Figure \ref{fig:data_collect}, our iterative pipeline details a methodology for extracting a large-scale corpus of mathematical content from Common Crawl, initiated from a small, high-quality set of mathematical data, termed the seed corpus. This iterative approach is generalizable and can be applied to other specialized domains such as computer code.

As an initial step, OpenWebMath \citep{openwebmath}, a curated compilation of high-quality mathematics web content, is selected to serve as the seed corpus. Using this seed, we develop a fastText model \citep{joulin2016fasttext} to retrieve additional web pages resembling the characteristics of OpenWebMath. For this, 500,000 samples are drawn at random from the seed corpus as positive instances, while another 500,000 web pages from Common Crawl are taken as negative instances. Model training employs an open-source implementation\footnote{\url{https://fasttext.cc}}, with hyperparameters set as follows: a vector size of 256, a learning rate of 0.1, a word n-gram length of up to 3, a minimum occurrence threshold of 3, and 3 training epochs. To address redundancy and reduce the Common Crawl’s scale, we use both URL-level deduplication and near-deduplication, yielding a corpus of 40 billion unique HTML documents. The trained fastText model is then applied to identify and recall web pages with mathematical relevance from this deduplicated subset. To ensure quality, pages are ranked using fastText-predicted scores and only the top-scoring pages are retained. We empirically determine the amount of data to preserve by pre-training models with the top 40B, 80B, 120B, and 160B tokens; for the initial iteration, we select the top 40B tokens for retention.

Following the initial round of data collection, a substantial portion of mathematical web pages remained uncaptured, primarily due to the insufficient heterogeneity in the set of positive examples used to train the fastText model. To address this limitation, we broaden the diversity of the seed corpus by incorporating additional sources of mathematical web content, thereby enhancing the fastText model’s coverage. Our process begins by partitioning the complete Common Crawl dataset into non-overlapping domains, with each domain encompassing web pages under the same base URL. For every domain, we compute the proportion of pages that were acquired during the first iteration. Domains in which more than 10\% of the pages are collected are designated as math-related domains (e.g., \url{mathoverflow.net}).

Next, we manually examine these domains to identify URLs that specifically denote mathematical content (e.g., \url{mathoverflow.net/questions}). We then augment the seed corpus by including any uncollected web pages associated with these annotated URLs. This enrichment process enables us to assemble a more comprehensive set of positive samples, facilitating the training of a more robust fastText model with greater recall for mathematical material in subsequent rounds. After performing four cycles of data collection, we compile a total of 35.5M mathematical web pages, which collectively comprise 120B tokens. Notably, by the fourth round, we observe that approximately 98\% of all available data was already retrieved by the third iteration, prompting us to conclude the data collection process at this point.

In order to prevent overlap with benchmark datasets, we adopt the methodology from \cite{deepseek-coder} to eliminate web pages containing content from English mathematical benchmarks such as GSM8K \citep{gsm8k} and MATH \citep{MATH}, as well as Chinese benchmarks like CMATH \citep{wei2023cmath} and AGIEval \citep{agieval}. The filtering protocol is as follows: any text fragment containing a 10-gram that exactly matches any substring within these evaluation benchmarks is purged from the math training dataset. For benchmark fragments between 3 and 10 grams in length, exact string matching is similarly applied to remove contaminated web content.

\subsection{Validating the Quality of the DeepSeekMath~Corpus}

To assess the relative quality of the DeepSeekMath~Corpus, we conduct pre-training experiments in comparison to other contemporary mathematical corpora:

\begin{itemize}[topsep=0pt]
    \item \textbf{MathPile} \citep{mathpile}: A composite corpus totaling 8.9B tokens, integrating material from textbooks, Wikipedia, ProofWiki, CommonCrawl, StackExchange, and arXiv, with the majority (exceeding 85\%) derived from arXiv;
    \item \textbf{OpenWebMath} \citep{openwebmath}: Mathematical content extracted from CommonCrawl, containing 13.6B tokens in total;
    \item \textbf{Proof-Pile-2} \citep{llemma}: A dataset incorporating OpenWebMath, AlgebraicStack (10.3B tokens of mathematical code), and arXiv documents (28.0B tokens). In our experiments with Proof-Pile-2, we adhere to the ratio of arXiv:Web:Code as 2:4:1, as suggested in \cite{llemma}.
\end{itemize}

\subsubsection{Training Setting}
\label{sec:quality-policy}
We conduct math-specific training on a broadly pre-trained language model with 1.3B parameters, adhering to the same architecture as the DeepSeek LLMs \citep{deepseek-llm}, and refer to this model as DeepSeek-LLM 1.3B. Each mathematical corpus is utilized to train a separate model for a total of 150B tokens. All training experiments leverage the resource-efficient HAI-LLM framework \citep{haillm}. Mirroring the training methodology established for DeepSeek LLMs, we employ the AdamW optimizer \citep{adamW} with hyperparameters $\beta_1=0.9$, $\beta_2=0.95$, and $\mathrm{weight\_decay}=0.1$. The learning rate follows a staged schedule: it linearly increases to its peak after 2,000 warmup steps, then decays to 31.6\% of its maximum by the 80\% mark of total training, and is further reduced to 10.0\% at the 90\% point. The peak learning rate is set at 5.3e-4. For each step, we use a batch size of 4M tokens and set the context window to 4K tokens.

\subsubsection{Evaluation Results}

\textbf{DeepSeekMath~Corpus distinguishes itself by superior quality, extensive multilingual coverage, and unprecedented scale.}

\begin{itemize}[topsep=0pt]
    \item \textbf{High-quality}: 
    The downstream capabilities of the models are assessed using eight different mathematical benchmarks with few-shot chain-of-thought prompting \cite{cot}. As illustrated in Table \ref{tab:corpora_comparison}, the model fine-tuned on the DeepSeekMath~Corpus consistently outperforms counterparts trained on alternative corpora. Further, Figure \ref{fig:corpora_comparisons} highlights that at the 50B token mark (representing one complete pass through Proof-Pile-2), the DeepSeekMath-trained model demonstrates stronger performance, evidencing a higher average quality for the DeepSeekMath~Corpus.
    \item \textbf{Multilingual}: 
    The DeepSeekMath~Corpus includes materials spanning a range of languages, with English and Chinese comprising the majority. Table \ref{tab:corpora_comparison} demonstrates that leveraging the DeepSeekMath~Corpus yields significant gains in mathematical reasoning for both English and Chinese datasets. In contrast, traditional math corpora, which are predominantly English-based, offer minimal benefit and in some cases even degrade performance on Chinese math tasks.
    \item \textbf{Large-scale}: 
    DeepSeekMath~Corpus vastly surpasses existing math corpora in magnitude. As shown in Figure \ref{fig:corpora_comparisons}, DeepSeek-LLM 1.3B models trained on DeepSeekMath~Corpus exhibit sharper improvements and maintain advancement throughout training. By comparison, baseline corpora are considerably smaller, necessitating multiple repeated passes during training and causing rapid stagnation in model performance.
\end{itemize}

\subsection{Training and Evaluating DeepSeekMath-Base 7B}

This section details the development of DeepSeekMath-Base 7B, a foundational model distinguished by its advanced mathematical reasoning proficiency. The initialization of our model is based on DeepSeek-Coder-Base-v1.5 7B \citep{deepseek-coder}, followed by a training regime spanning 500B tokens. The training data composition comprises 56\% DeepSeekMath Corpus, 4\% AlgebraicStack, 10\% arXiv, 20\% Github code, and the final 10\% consists of Common Crawl natural language data in both English and Chinese. Training generally follows the protocol outlined in Section \ref{sec:quality-policy}, with two exceptions: the learning rate is capped at 4.2e-4, and the batch size is increased to 10M tokens.

To rigorously characterize the mathematical aptitude of DeepSeekMath-Base 7B, we perform extensive evaluations that test its competence in constructing self-sufficient mathematical answers, leveraging computational tools for problem solving, and engaging in formal theorem proving. We supplement these mathematical assessments by presenting a broader overview of the model’s general abilities, including its performance in natural language comprehension, logical reasoning, and code generation.

\paragraph{Step-by-Step Reasoning for Mathematical Problem Solving}

DeepSeekMath-Base’s aptitude for mathematical problem-solving is benchmarked through few-shot chain-of-thought prompting \citep{cot} on eight datasets spanning both English and Chinese. These benchmarks include tasks requiring quantitative reasoning, such as GSM8K \citep{gsm8k}, MATH \citep{MATH}, and CMATH \citep{wei2023cmath}, as well as multiple-choice evaluations like MMLU-STEM \citep{mmlu} and Gaokao-MathQA \citep{agieval}. The selection of benchmarks ensures coverage of a wide range of mathematical domains, addressing content from elementary school through collegiate-level material.

As indicated in Table \ref{tab:base_math_cot}, among the open-source base models—including prominent options such as Mistral 7B \citep{mistral} and the newly introduced Llemma 34B \citep{llemma} trained with Proof-Pile-2 \citep{llemma}—DeepSeekMath-Base 7B achieves the highest scores across all eight assessed benchmarks. Remarkably, on the challenging MATH dataset that targets competition-level problems, DeepSeekMath-Base achieves an absolute improvement exceeding 10\% over other open-source base models. It also surpasses Minerva 540B \citep{minerva}, a proprietary model based on PaLM \citep{palm} and extended with mathematical corpus, despite Minerva being 77 times larger in scale.

\paragraph{Mathematical Problem Solving with Tool Use} 
We assess models' capabilities in program-assisted mathematical problem solving using GSM8K and MATH, employing few-shot program-of-thought prompting as described in \citep{pot,pal}. Here, the models are prompted to address each task by crafting Python scripts, leveraging libraries like \textit{math} and \textit{sympy} for handling complex calculations. The script’s output is then taken as the predicted answer. Table \ref{tab:base_math_pot_proof} shows that DeepSeekMath-Base 7B establishes a new state of the art, outperforming the previous best, Llemma 34B.

\paragraph{Formal Mathematics}

Automating the generation of formal proofs is essential for improving the rigor, accuracy, and efficiency of mathematical proof construction and has garnered growing interest. We benchmark DeepSeekMath-Base 7B on informal-to-formal proving tasks outlined in \citep{dsp_proof}, where the objective is to translate an informal conjecture, its formalized version, and an informal argument into a formal proof. The miniF2F dataset \citep{minif2f}—which features Olympiad-style formal mathematics—is used, and few-shot prompting is applied to produce Isabelle-formatted proofs for each item. Following \cite{dsp_proof}, we employ model-generated proof outlines and utilize the Sledgehammer \citep{sledgehammer} automated prover to elaborate missing proof steps. As reported in Table \ref{tab:base_math_pot_proof}, DeepSeekMath-Base 7B delivers robust outcomes in the task of proof autoformalization.

\paragraph{Natural Language Understanding, Reasoning, and Code}

To measure performance on broader cognitive tasks, we evaluate the models on MMLU \citep{mmlu} for language comprehension, BBH \citep{bbh} for reasoning, and HumanEval \citep{codex} as well as MBPP \citep{mbpp} for coding abilities. Results in Table \ref{tab:understanding_reasoning_code} reveal that DeepSeekMath-Base 7B attains marked gains on MMLU and BBH relative to its predecessor, DeepSeek-Coder-Base-v1.5 \citep{deepseek-coder}, highlighting the advantage of mathematical training on language understanding and reasoning tasks. Furthermore, with continued training on code data, DeepSeekMath-Base 7B preserves the coding performance exhibited by DeepSeek-Coder-Base-v1.5 on HumanEval and MBPP. Across all three reasoning and coding benchmarks, DeepSeekMath-Base 7B consistently outperforms the general-purpose Mistral 7B model \citep{mistral}.

\section{Supervised Fine-Tuning}

\subsection{SFT Data Curation}

We develop an instruction-tuning dataset for mathematics that encompasses both English and Chinese questions, representing a broad range of mathematical domains and varying difficulty: each problem is accompanied by solutions employing chain-of-thought (CoT) \citep{cot}, program-of-thought (PoT) \citep{pot,pal}, or tool-integrated reasoning formats \citep{tora}. The overall training set comprises 776K examples.

\begin{itemize}[topsep=0pt]
    \item \textbf{English mathematical datasets}: Our English resources consist of GSM8K and MATH items annotated with tool-integrated responses, supplemented with select data from MathInstruct \citep{MathInstruct} and the Lila-OOD training split \citep{lila} where solutions leverage either CoT or PoT. The dataset encompasses multiple mathematical subfields, such as algebra, calculus, probability, geometry, and number theory.
    \item \textbf{Chinese mathematical datasets}: For Chinese, we curate K-12 level mathematical items distributed across 76 sub-categories—including, for instance, linear equations—with answer rationales prepared in both CoT and tool-integrated formats.
\end{itemize}

\subsection{Training and Evaluating DeepSeekMath-Instruct 7B}

This section details the procedure for instruction-tuning DeepSeekMath-Instruct 7B, which is initialized from DeepSeekMath-Base. During training, samples are concatenated in random order until the combined input reaches the 4K token limit per instance. Training proceeds for 500 steps with a constant learning rate of 5e-5 and a batch size of 256.

Mathematical capabilities are assessed with and without external tool integration across four quantitative reasoning benchmarks in both English and Chinese. We conduct comparative evaluations against prominent contemporaneous models:

\begin{itemize}[topsep=0pt]
    \item \textbf{Closed-source models} considered are: (1) the GPT suite—especially GPT-4 \citep{gpt4} and GPT-4 Code Interpreter~\footnote{\url{https://openai.com/blog/chatgpt-plugins\#\#code-interpreter}}, (2) Gemini Ultra and Pro \citep{gemini}, (3) Inflection-2 \citep{inflection-2}, (4) Grok-1~\footnote{\url{https://x.ai/model-card}}, as well as competitive models from Chinese developers including (5) Baichuan-3~\footnote{\url{https://www.baichuan-ai.com}} and (6) the latest GLM-4~\footnote{\url{https://open.bigmodel.cn/dev/api\#glm-4}}
    from the GLM series \citep{glm}.
\end{itemize}

These models are generally designed for broad applications and most have been subjected to alignment techniques.

\item \textbf{Open-source models} span: standard-purpose LLMs such as (1) DeepSeek-LLM-Chat 67B \citep{deepseek-llm}, (2) Qwen 72B \citep{qwen}, (3) SeaLLM-v2 7B \citep{seallm}, and (4) ChatGLM3 6B \citep{chatglm3}. Models specialized for mathematical tasks include (5) InternLM2-Math 20B~\footnote{\url{https://github.com/InternLM/InternLM-Math}}, which extends InternLM2 and adds instruction tuning following math-centered pretraining; (6) Math-Shepherd-Mistral 7B, incorporating PPO optimization \citep{schulman2017proximal} into Mistral 7B \citep{mistral} using a process-supervised reward function; (7) the WizardMath series \citep{wizardmath}, which advances mathematical reasoning in both Mistral 7B and Llama-2 70B \citep{llama2} via evolve-instruct (an approach that leverages AI-generated instructions for instruction tuning) and PPO, with training primarily drawing from GSM8K and MATH; (8) MetaMath 70B \citep{metamath}, consisting of Llama-2 70B refined on an expanded version of GSM8K and MATH; (9) ToRA 34B \cite{tora}, representing CodeLlama 34B tailored for mathematical reasoning with integrated tool use; (10) MAmmoTH 70B \citep{MathInstruct}, which results from MathInstruct-based instruction tuning of Llama-2 70B.

As reported in Table \ref{tab:sft_rl_math}, in settings that do not permit the use of external tools, DeepSeekMath-Instruct 7B exhibits robust stepwise reasoning abilities. Notably, on the challenging MATH benchmark, our model outperforms all other open-source systems and most proprietary models (including Inflection-2 and Gemini Pro) by a margin of at least 9 percentage points. This holds true even against much larger models (e.g., Qwen 72B) and those explicitly refined using mathematical reinforcement learning, such as WizardMath-v1.1 7B. Although DeepSeekMath-Instruct achieves performance comparable to Chinese proprietary models like GLM-4 and Baichuan-3 on the MATH dataset, it remains behind GPT-4 and Gemini Ultra.

In settings where models may combine natural language reasoning with the use of programmatic tools, DeepSeekMath-Instruct 7B attains nearly 60\% accuracy on MATH, establishing a new high among open-source models. On other assessment tasks, it remains competitive with DeepSeek-LLM-Chat 67B, which was previously state-of-the-art but operates at ten times the parameter scale.

\section{Reinforcement Learning}

\subsection{Group Relative Policy Optimization}
The use of reinforcement learning (RL) has shown considerable efficacy for boosting LLMs' mathematical reasoning skills following Supervised Fine-Tuning (SFT) \citep{wang2023math, wizardmath}. This section presents our proposed RL strategy—Group Relative Policy Optimization (GRPO)—which balances efficiency and effectiveness.

\subsubsection{From PPO to GRPO}
Proximal Policy Optimization (PPO) \citep{schulman2017proximal} stands out as a mainstream actor-critic RL method adopted in LLM RL fine-tuning \citep{ouyang2022training}. PPO seeks to update model parameters by maximizing the following objective:

\begin{equation}
\footnotesize
    \mathcal{J}_{PPO}(\theta) = \mathbb{E}{[q \sim P(Q), o \sim \pi_{\theta_{old}}(O|q)]} \frac{1}{|o|} \sum_{t=1}^{|o|} \min \left[ \frac{\pi_\theta(o_{t} | q, o_{<t})}{\pi_{\theta_{old}}(o_{t} | q, o_{<t})} A_{t}, \text{clip} \left( \frac{\pi_\theta(o_{t} | q, o_{<t})}{\pi_{\theta_{old}}(o_{t} | q, o_{<t})}, 1 - \varepsilon, 1 + \varepsilon \right)  A_{t} \right] ,
\end{equation}

Here, $\pi_{\theta}$ denotes the current policy, while $\pi_{\theta_{old}}$ stands for the previous iteration of the policy. The variables $q$ and $o$ refer to the questions and generated outputs, respectively, which are drawn from the question dataset and the historical policy $\pi_{\theta_{old}}$. The parameter $\varepsilon$ serves as a clipping hyper-parameter in PPO, ensuring stable optimization. The term $A_t$ represents the advantage estimate, typically obtained using Generalized Advantage Estimation (GAE) \citep{gae}, which relies on the set of future rewards $\{r_{\ge t}\}$ and a learned value function $V_{\psi}$. Consequently, PPO necessitates concurrent training of a value function alongside the policy model. To avoid excessive reward model exploitation, a standard technique is to incorporate a token-level KL divergence penalty against a reference model into the reward at every token position, as introduced by \citet{ouyang2022training}, giving rise to

\begin{equation}
    r_{t} = r_\varphi(q, o_{\le t}) - \beta \log\frac{\pi_{\theta}(o_{t}|q, o_{<t})}{\pi_{ref}(o_{t}|q, o_{<t})},
\label{eq:PPO-reward}
\end{equation}

where $r_\varphi$ denotes the reward model, $\pi_{ref}$ is typically set as the initial SFT model serving as the reference, and $\beta$ quantifies the influence of the KL penalty.

Due to the fact that the value function used in PPO is often a separate network of similar scale to the policy model, this approach significantly increases both computational and memory demands. Moreover, during reinforcement learning, the value function is included as a baseline to reduce the variance of the advantage estimates. In large language models, it is common that only the final token is awarded a reward by the reward model, complicating efforts to train a reliable value function for every token. To overcome these challenges, we introduce Group Relative Policy Optimization (GRPO) (as illustrated in Figure \ref{fig:grpo}), which eliminates the requirement for an auxiliary value function as seen in PPO. Instead, GRPO uses the mean reward across a set of sampled responses to the same question as a baseline. In particular, for each question $q$, a group of responses $\{o_1, o_2, \cdots, o_G\}$ is sampled from $\pi_{\theta_{old}}$, and the objective optimized is as follows:

\begin{equation}
\footnotesize
\begin{split}
    \mathcal{J}_{GRPO}(\theta) &= \mathbb{E}{[q \sim P(Q), \{o_i\}_{i=1}^G \sim \pi_{\theta_{old}}(O|q)]}  \\
    & \frac{1}{G}\sum_{i=1}^G\frac{1}{|o_i|} \sum_{t=1}^{|o_i|} \left\{ \min \left[ \frac{\pi_\theta(o_{i,t} | q, o_{i,<t})}{\pi_{\theta_{old}}(o_{i,t} | q, o_{i,<t})} \hat{A}_{i,t}, \text{clip} \left( \frac{\pi_\theta(o_{i,t} | q, o_{i,<t})}{\pi_{\theta_{old}}(o_{i,t} | q, o_{i,<t})}, 1 - \varepsilon, 1 + \varepsilon \right)  \hat{A}_{i,t} \right] - \beta \mathbb{D}_{KL}\left[\pi_{\theta} || \pi_{ref}\right]\right\} ,
\end{split}
\label{eq:GRPO-obj}
\end{equation}

where $\varepsilon$ and $\beta$ are both hyper-parameters, and $\hat{A}_{i,t}$ is an advantage value calculated strictly from the relative rewards among outputs within each sampled group, with further discussion provided in the subsequent sections. This group-relative approach for advantage computation in GRPO is particularly congruent with the comparative design of reward models, as such models are frequently trained using data that contrasts responses to the same prompt. It is worth noting that GRPO incorporates the KL regularization term directly in the loss function—rather than inside the reward as in PPO—thereby simplifying the estimation of $\hat{A}_{i,t}$. Furthermore, distinct from the reward-based KL penalty described in (\ref{eq:PPO-reward}), we compute the KL divergence utilizing the following unbiased

\begin{algorithm}[t]
  \small
  \caption{Iterative Group Relative Policy Optimization}
  \textbf{Input}: Initial policy model $\pi_{\theta_{\text{init}}}$; reward models $r_\varphi$; task prompt set $\mathcal{D}$; 
  hyperparameters $\varepsilon$, $\beta$, $\mu$
  \begin{algorithmic}[1]
    \State Initialize policy model: $\pi_\theta \leftarrow \pi_{\theta_{\text{init}}}$
    \For{iteration = 1, \dots, I}
       \State Set reference model $\pi_{ref} \leftarrow \pi_{\theta}$
      \For{step = 1, \dots, M}
      \State Draw a batch $\mathcal{D}_b$ from $\mathcal{D}$
      \State Update old policy model $\pi_{\theta_{old}} \leftarrow \pi_{\theta}$
      \State For each prompt $q \in \mathcal{D}_b$, generate $G$ outputs $\{o_i\}_{i=1}^G$ by sampling from $\pi_{\theta_{old}}(\cdot \mid q)$
      \State Evaluate each generated $o_i$ to obtain corresponding rewards $\{r_i\}_{i=1}^G$ via $r_{\varphi}$
      \State Use group relative advantage estimation to calculate $\hat{A}_{i,t}$ for the $t$-th token in each $o_i$
      \For{GRPO iteration = 1, \dots, $\mu$}
        \State Update policy parameters of $\pi_{\theta}$ to maximize the GRPO objective (Equation \ref{eq:GC-GRPO})
      \EndFor
    \EndFor
    \State Continuously train $r_\varphi$ using a replay strategy.
    \EndFor
  \end{algorithmic}
  \textbf{Output}: $\pi_\theta$
  \label{alg:iter-grpo}
\end{algorithm}

\subsubsection{Outcome Supervision RL with GRPO} Consider, for each question $q$, the procedure where a collection of outputs $\{o_1, o_2, \cdots, o_G\}$ is sampled from the current policy snapshot $\pi_{\theta_{old}}$. A reward model assigns a score to each output, producing a vector of $G$ rewards $\mathbf{r} = \{r_1, r_2, \cdots, r_G\}$. These reward values are subsequently standardized: the mean of the group is subtracted and then the result is divided by the group standard deviation. In this outcome supervision approach, the normalized reward is assigned to the end of each output $o_i$, and this value is broadcast as the advantage $\hat{A}_{i, t}$ for all positions $t$ in the corresponding output; formally, $\hat{A}_{i, t} = \widetilde{r}_i = \frac{r_i- {\rm mean}(\mathbf{r})}{{\rm std}(\mathbf{r})}$. Policy optimization is then conducted by maximizing the objective specified in Equation (\ref{eq:GRPO-obj}).

\subsubsection{Process Supervision RL with GRPO} In outcome supervision, only a terminal reward is given for each sequence, which can be limiting for supervising intricate tasks such as mathematics. Building on prior work \cite{wang2023math}, we additionally implement process supervision, which involves assigning feedback at each intermediate reasoning stage. Specifically, for a given input $q$ and $G$ sampled outputs $\{o_1, o_2, \cdots, o_G\}$, a process reward model evaluates and scores each reasoning step, resulting in rewards structured as $\mathbf{R} = \{ \{r_1^{index(1)},\cdots,r_1^{index(K_1)}\}, \cdots,  \{r_G^{index(1)},\cdots,r_G^{index(K_G)}\} \}$, where $index(j)$ denotes the index of the terminal token for the $j$-th step in $o_i$, and $K_i$ represents the number of steps in output $o_i$. These stepwise rewards are normalized via subtraction of the mean and division by the standard deviation of $\mathbf{R}$: $\widetilde{r}_i^{index(j)} = \frac{r_i^{index(j)} - {\rm mean(\mathbf{R})}}{{\rm std(\mathbf{R})}}$. The advantage for each token $t$ within $o_i$ is then determined by summing all normalized future step rewards,

\subsubsection{Iterative RL with GRPO} During reinforcement learning, the previously trained reward model may become inadequate to effectively guide the evolving policy model. To address this, we investigate an iterative approach to RL utilizing GRPO. As detailed in Algorithm \ref{alg:iter-grpo}, the iterative GRPO framework involves constructing new reward model training datasets using samples generated from the latest policy model. The reward model is then further optimized by applying a replay buffer containing 10\% of its prior training data. Subsequently, we set the reference model to align with the current policy model and continue the optimization of the policy model using the updated reward model.

\subsection{Training and Evaluating DeepSeekMath-RL}

Our reinforcement learning experiments employ DeepSeekMath-Instruct 7B as the base model. The RL training corpus is composed of chain-of-thought-formulated problems from GSM8K and MATH, sourced from the SFT data and totaling approximately 144K questions. Other SFT items are omitted in order to analyze the influence of RL specifically on benchmarks without overlapping data during RL training. The construction of reward model training sets follows the methodology of \citep{wang2023math}. The initial reward model is fine-tuned on DeepSeekMath-Base 7B using a learning rate of 2e-5. For GRPO optimization, the policy model’s learning rate is set to 1e-6, with a KL regularization coefficient of 0.04. For each question, $64$ sample completions are produced. The maximum input length is limited to 1024 tokens, and batch size is set at 1024. Only one policy model update is conducted after each exploration phase. DeepSeekMath-RL 7B’s evaluation protocols mirror those for DeepSeekMath-Instruct 7B. Within this setup, GSM8K and MATH tasks featuring chain-of-thought solutions are classified as in-domain, whereas remaining evaluation benchmarks are treated as out-of-domain tasks.

Table \ref{tab:sft_rl_math} presents the results achieved by both open- and closed-source models employing chain-of-thought and tool-augmented reasoning techniques across English and Chinese benchmarks. Our observations include: 1) DeepSeekMath-RL 7B reaches 88.2\% and 51.7\% accuracy on GSM8K and MATH benchmarks, respectively, when utilizing chain-of-thought reasoning. These results outperform every open-source model in the 7B to 70B parameter range, as well as most closed-source counterparts. 2) Notably, DeepSeekMath-RL 7B's training is limited to chain-of-thought-formatted instruction tuning data from GSM8K and MATH, beginning with DeepSeekMath-Instruct 7B. Despite this restriction in training data diversity, DeepSeekMath-RL 7B consistently exceeds DeepSeekMath-Instruct 7B across all tested metrics, highlighting the strengths of reinforcement learning.

\section{Discussion}
This section discusses our insights derived from both pre-training and reinforcement learning experiments.

\subsection{Lessons Learnt in Pre-Training}

We begin by discussing our pre-training experience. Unless otherwise mentioned, the training configuration adheres to that specified in Section \ref{sec:quality-policy}. Additionally, references to the DeepSeekMath Corpus in this section correspond to an 89B-token dataset compiled during the second round of data collection.

\subsubsection{Code Training Benefits Mathematical Reasoning}

There is a widely held, though largely untested, belief that code training can enhance reasoning capabilities. We aim to provide empirical support for this notion in the mathematical context: code training improves models’ proficiency in mathematical reasoning tasks, irrespective of whether external tools are used.

To investigate the impact of code training on mathematical reasoning, we explored both two-stage and single-stage training setups:

\textbf{Two-Stage Training}

\begin{itemize}[topsep=0pt]
    \item \textbf{Code Training for 400B Tokens $\rightarrow$ Math Training for 150B Tokens}: The DeepSeek-LLM 1.3B model is first trained on 400B tokens sourced from code, followed by a secondary training phase with 150B tokens of mathematical content;
    \item \textbf{General Training for 400B Tokens $\rightarrow$ Math Training for 150B Tokens}: For comparison, we conduct an alternative training regimen where, instead of code, the model is initially exposed to 400B general tokens drawn from an expansive corpus curated by DeepSeek-AI. This is done to scrutinize whether pretraining on code as opposed to general content has measurable effects on the subsequent development of mathematical reasoning skills.
\end{itemize}

\textbf{One-Stage Training}

\begin{itemize}[topsep=0pt]
    \item \textbf{Math Training for 150B Tokens}: Here, DeepSeek-LLM 1.3B undergoes a single-stage training focused solely on 150B tokens of math-specific material;
    \item \textbf{Training on a mixture of 400B Code Tokens and 150B Math Tokens}: Since follow-up mathematical training post code training tends to deteriorate coding capabilities, we further test whether a joint, single-phase training blending code and mathematical data (in a combined 400B and 150B token regime) could preserve coding abilities while simultaneously supporting mathematical reasoning and preventing catastrophic forgetting.
\end{itemize}

\paragraph{Results}
As presented in Table \ref{tab:code-math} and Table \ref{tab:code-math-general-eval}, we summarize the downstream task performance yielded by these distinct training strategies.

In both sequential (two-stage) and combined (one-stage) training frameworks, prior exposure to code provides notable benefits to program-assisted mathematical problem solving. Specifically, the two-stage protocol reveals that code pretraining already confers a considerable boost to DeepSeek-LLM’s proficiency on tasks such as GSM8K and MATH when solutions are programmatically expressed in Python, with a further performance increment observed after math-specific fine-tuning. Notably, simultaneous training with code and math tokens, in a one-stage setup, appears to significantly counteract the loss of coding ability typically seen with two-stage approaches and delivers synergetic improvements for both code-centric tasks (see Table \ref{tab:code-math-general-eval}) and for program-assisted math reasoning (see Table \ref{tab:code-math}).

Beyond tool-augmented reasoning, code-centric pretraining is also shown to moderately enhance direct mathematical reasoning skills without external assistance during the initial training stage, with additional math finetuning resulting in peak accuracy. However, the mixed-modality one-stage strategy, wherein code and math tokens are combined, tends to undermine such unaided mathematical reasoning, possibly due to capacity constraints of the DeepSeek-LLM 1.3B model, which may struggle to internalize both code and mathematics effectively when presented concurrently.

\subsubsection{ArXiv Papers Appear Unhelpful for Enhancing Mathematical Reasoning}

Inclusion of ArXiv papers is standard practice in math pre-training datasets \citep{minerva,gpt-f,llemma,mathpile}, yet comprehensive examinations of their actual effect on mathematical reasoning remain limited. Somewhat unexpectedly, our empirical results indicate that ArXiv papers do not substantially enhance mathematical reasoning abilities. To investigate this, we evaluate several model scales, namely DeepSeek-LLM 1.3B and DeepSeek-Coder-Base-v1.5 7B \citep{deepseek-coder}, using ArXiv datasets subject to distinct preprocessing approaches:

\begin{itemize}[topsep=0pt]
    \item \textbf{MathPile} \citep{mathpile}: an 8.9B-token collection processed with heuristic cleaning and filtering procedures, comprising more than 85\% scientific ArXiv documents;
    \item \textbf{ArXiv-RedPajama} \citep{redpajama}: the complete set of ArXiv LaTeX submissions with extraneous content (such as preambles, comments, macros, and bibliographies) stripped out, totaling 28.0B tokens.
\end{itemize}

For our analysis, we separately pre-train DeepSeek-LLM 1.3B on 150B tokens and DeepSeek-Coder-Base-v1.5 7B on 40B tokens from each ArXiv dataset. The experimental findings suggest that ArXiv materials provide minimal, if any, benefit for mathematical reasoning: when models are trained solely on ArXiv data, they show no significant gains—and sometimes even worse results—across several mathematical evaluation suites of varying challenge levels considered here. These test sets comprise quantitative reasoning resources (such as GSM8K and MATH, see Table \ref{tab:arxiv-cot}), multiple-choice tests (including MMLU-STEM, Table \ref{tab:arxiv-cot}), and formal mathematics assessments (like miniF2F, Table \ref{tab:arxiv_atp}).

Nonetheless, this observation is not without caveats. Certain aspects remain unexamined:

\begin{itemize}[topsep=0pt]
    \item The influence of ArXiv-derived data on niche mathematical skills not addressed in this work, such as translating formal mathematical content into informal descriptions;
    \item The interplay between ArXiv papers and other dataset sources in composite training regimes;
    \item The possibility that advantageous effects of ArXiv data could emerge with increased model capacity.
\end{itemize}

Consequently, more thorough investigation is warranted, which we defer to future research.

\subsection{Reflections on Reinforcement Learning}
\subsubsection{In Pursuit of a Cohesive Framework} Here, we outline a cohesive framework for analyzing diverse training algorithms, encompassing SFT, RFT, DPO, PPO, GRPO, and extend our work with empirical studies to identify influential elements within this framework. Generally, the parameter gradient for these training schemes with respect to $\theta$ can be described as:

\begin{equation}
    \nabla_{\theta}\mathcal{J}_{\textcolor{red}{\mathcal{A}}}(\theta) = \mathbb{E}[\underbrace{(q,o) \sim \textcolor{red}{\mathcal{D}}}_{Data \ Source}]\left( \frac{1}{|o|} \sum_{t=1}^{|o|}  \underbrace{GC_{{\mathcal{A}}}(q, o, t, \textcolor{red}{\pi_{{rf}}})}_{Gradient \ Coefficient}  \nabla_{\theta}\log \pi_{\theta}(o_t | q, o_{<t})\right).
\label{eq:objective}
\end{equation}

This formulation incorporates three fundamental elements: 1) \textit{Data Source $\mathcal{D}$}, specifying the corpus for training; 2) \textit{Reward Function $\pi_{{rf}}$}, which furnishes the reward signals that guide optimization; 3) \textit{Algorithm $\mathcal{A}$}, which manipulates both training inputs and rewards to yield the gradient coefficient $GC$, thereby modulating the intensity of penalization or reward assigned to each instance. Building upon this general framework, we examine several widely used approaches as follows:

\begin{itemize}
    \item  \textbf{Supervised Fine-tuning (SFT)}: In this setup, the model undergoes fine-tuning using SFT data curated by human annotators.
    \item \textbf{Rejection Sampling Fine-tuning (RFT)}: Here, the SFT model is further adapted by fine-tuning it on a filtered subset of outputs, sampled from itself, where filtering is performed based on answer correctness for the SFT-provided questions.
    \item \textbf{Direct Preference Optimization (DPO)}: DPO enhances the SFT model by employing additional outputs generated via the SFT model and optimizing a pairwise DPO objective.
    \item \textbf{Online Rejection Sampling Fine-tuning (Online RFT)}: In contrast to standard RFT, Online RFT begins with the SFT-initialized policy and performs ongoing fine-tuning using outputs sampled from the actively evolving policy model.
    \item \textbf{PPO/GRPO}: With PPO/GRPO, the policy model is initialized from the SFT checkpoint and trained further by reinforcement, utilizing outputs sampled from the current (online) policy itself.
\end{itemize}

A summary of the methodological components is provided in Table \ref{tab:data_policy}. For an expanded discussion of the derivation steps, consult Appendix \ref{app:analysis-rl}.

\paragraph{Observation about Data Source} We categorize data sources into online sampling and offline sampling. In online sampling, training data originate from trajectories generated by the policy model being optimized during training. Conversely, offline sampling utilizes training data derived from outputs of the initial SFT model. Among the considered methods, RFT and DPO employ the offline approach, whereas Online RFT and GRPO leverage online sampling.

Figure \ref{fig:rl-analysis} illustrates that Online RFT substantially outperforms RFT across both evaluation benchmarks. At the early stages of optimization, Online RFT exhibits performance akin to RFT; however, as training progresses, Online RFT establishes a pronounced lead, reflecting the benefits of the online approach. This outcome aligns with expectations: initially, the actor and SFT model are very similar, resulting in sampled data with negligible divergence. Over time, as the actor evolves, its samples differ more markedly, allowing online data collection to deliver greater performance improvements.

\paragraph{Observation about Gradient Coefficient} The process by which each algorithm transforms reward signals into gradient coefficients, which guide model parameter updates, is a central distinction. Our experiments consider two reward function paradigms: `Rule,' where response quality is assessed according to the correctness of the output, and `Model,' where a trained reward model scores each response using rule-based annotations as ground truth. A salient distinction highlighted in Equations \ref{eq:GC-RFT} and \ref{eq:GC-GRPO} is that GRPO dynamically adapts its gradient coefficient in accordance with the reward model's output, permitting nuanced scaling of reinforcement and penalties based on reward magnitude. By contrast, Online RFT does not incorporate this mechanism—it applies a uniform level of reinforcement to all correctly answered responses and does not assign penalties to incorrect outputs.

As illustrated in Figure \ref{fig:rl-analysis}, GRPO outperforms the online RFT method, underscoring the effectiveness of tuning both positive and negative gradient coefficients. Moreover, GRPO+PS yields better results than GRPO+OS, demonstrating the advantages of employing detailed, step-sensitive gradient adjustments. Additionally, we investigate the impact of iterative RL by executing two iterations in our trials. The results in Figure \ref{fig:iter-rl} reveal that introducing iterations in RL markedly enhances performance, with the most notable gains observed after the initial iteration.

\subsubsection{Why RL Works?} In this study, we apply reinforcement learning using a subset derived from instruction tuning data, resulting in marked improvements over the baseline instruction-tuned model. To provide further insight into the effectiveness of reinforcement learning, we assess both Pass@K and Maj@K accuracies for the Instruct and RL models across two benchmark datasets. According to Figure \ref{fig:maj-pass}, RL improves Maj@K scores but has little effect on Pass@K. This suggests that reinforcement learning bolsters overall performance by making the distribution of model outputs more resilient, specifically by enhancing the likelihood that a correct answer appears among the TopK options, rather than by fundamentally upgrading the model’s core abilities. Comparable observations were made by \citep{wang2023making}, who identified a \textbf{misalignment problem} in reasoning tasks handled by SFT models, noting that various preference alignment methods \citep{yuan2023rrhf,song2023preference,wang2023making} can notably strengthen the reasoning capacity of these models.

\subsubsection{How to Achieve More Effective RL?}
\label{sec:effec_rl}
Our experiments confirm the strong efficacy of reinforcement learning on mathematical reasoning tasks. Furthermore, we propose a unified framework that conceptualizes diverse representative training strategies as variants or simplifications of RL-based techniques. Within this structure, as presented in Equation \ref{eq:objective}, three essential components emerge: Data Source, Algorithm, and Reward Function. We outline several promising future research avenues concerning each of these elements.

\paragraph{Data Source} The data source forms the foundation of any training approach. In reinforcement learning scenarios, we define the data source as consisting of unlabeled queries paired with outputs generated from the policy model. In our current work, we restrict ourselves to prompts originating from the instruction tuning phase, applying straightforward nucleus sampling to obtain outputs. We hypothesize that this limitation could explain why our RL framework exclusively elevates Maj@K performance. For future studies, we intend to broaden our RL approach to encompass out-of-distribution queries and incorporate \textbf{advanced sampling (decoding) strategies} such as those using tree-search methodologies \citep{yao2023tree}. Furthermore, we highlight that \textbf{efficient inference techniques} \citep{xia-etal-2023-speculative,leviathan2023fast,kwon2023efficient,xia2024unlocking}, which directly influence the exploratory capabilities of policy models, are also critical for maximizing effectiveness.

\paragraph{Algorithms} Algorithms utilize both the observed data and the reward signal to adjust the gradient coefficient, thereby modifying the model parameters. As indicated in Equation \ref{eq:objective}, existing approaches fundamentally place full \textbf{TRUST} in the reward function's feedback to either promote or diminish the likelihood of generating particular tokens. Nevertheless, there is no guarantee that the reward signal is always dependable, particularly when dealing with highly intricate tasks. For instance, the PRM800K dataset \citep{lightman2023let}, despite meticulous labeling efforts by skilled annotators, still exhibits around 20\% erroneous annotations\footnote{\url{https://github.com/openai/prm800k/issues/12\#issuecomment-1728491852}}. To address this, our focus shifts toward reinforcement learning algorithms capable of maintaining robustness in the presence of noisy reward signals. We contend that adopting \textbf{WEAK-TO-STRONG} \citep{burns2023weak} alignment strategies could profoundly transform current learning paradigms.

\paragraph{Reward Function} The reward function serves as the fundamental provider of learning signals during training. In reinforcement learning scenarios, the reward function commonly takes the form of a neural reward model. We identify three pivotal research trajectories regarding reward models: 1) \textbf{Improving the generalization capacity of reward models.} For reinforcement learning to meaningfully enhance the abilities of LLMs, reward models must generalize effectively to out-of-distribution prompts and sophisticated generative outputs; otherwise, RL may merely preserve the status quo in LLM output distributions; 2) \textbf{Capturing and representing reward model uncertainty.} Quantifying uncertainty has the potential to serve as a conceptual connection between unreliable reward signals and weak-to-strong alignment approaches; 3) \textbf{Constructing high-quality, process-level reward models efficiently} to supply detailed training feedback throughout complex reasoning chains \citep{lightman2023let,wang2023math}.

\section{Conclusion, Limitation, and Future Work}

In this work, we introduce DeepSeekMath, which establishes a new state-of-the-art among open-source models on the challenging MATH benchmark and delivers performance close to that of proprietary systems. DeepSeekMath~originates from DeepSeek-Coder-v1.5 7B, followed by an extended training regimen consisting of 500B tokens, notably including 120B math-focused tokens obtained from Common Crawl. Our comprehensive ablation experiments reveal that web-sourced documents are an underutilized yet valuable resource for high-quality mathematical data, while data from arXiv appears to have limited marginal benefit. We also propose Group Relative Policy Optimization (GRPO), building on Proximal Policy Optimization (PPO), which effectively enhances mathematical reasoning skills while being more memory-efficient. Empirical evaluations demonstrate that GRPO is effective, even for DeepSeekMath-Instruct 7B at already elevated benchmark performance levels. Additionally, we provide a unified framework for interpreting multiple methods and highlight several avenues for advancing reinforcement learning approaches.

Despite DeepSeekMath's~strong results on quantitative reasoning benchmarks, its performance on geometry and theorem-proving tasks remains comparatively limited when measured against closed-source models. For example, preliminary evaluations indicate the model struggles with geometric questions, such as those involving triangles and ellipses, suggesting possible data selection imbalances during both pre-training and fine-tuning phases. Furthermore, due to model size constraints, DeepSeekMath's~few-shot learning ability lags behind GPT-4. While GPT-4 exhibits noticeable gains when given few-shot inputs, DeepSeekMath~shows little difference between zero-shot and few-shot scenarios. Moving forward, we plan to refine our data selection pipeline further to create an even higher-quality pre-training corpus. We also intend to pursue research into more effective reinforcement learning techniques for large language models, as discussed in Section \ref{sec:effec_rl}.

\end{document}